% !TeX root = ../../main.tex
\subsection{Exact recursions and their integration constants}
\label{subsec:s07:03}

We first recover enough of the free energy to justify integrating the
field equation. The string equation removes the zero-degree time, and
the dilaton equation then sums the degree-one time. For fixed positive
defect, this leaves finitely many higher times. Their expression in
inverse-function coordinates will also give the condition at the
infinite end of a characteristic.

For integers $d\geq1$ and $k\geq1$, define the shifted times and
inverse-function jets by
\begin{equation}\label{eq:s07:jets}
 \mathcal I_d=\sum_{a\geq0}t_{d+a}\frac{v^a}{a!},\qquad
 W_k=W^{(k)}(v),\qquad
 W_1=1-\mathcal I_1,\quad W_k=-\mathcal I_k\ (k\geq2).
\end{equation}
The symbol $W$ here is the full-time series in
\eqref{eq:s07:genus-zero-inverse}; after specialization it is
$w_\theta$. A rational expression in the jets means an expression in
finitely many $W_k$, with denominators powers of $W_1$.

\begin{lmm}[Finite jet expressions]\label{lmm:s07:finite-jets}
For $g\geq2$ and $1\leq h\leq g$,
\begin{equation}\label{eq:s07:finite-jets}
 \mathcal F_{g,h}
 =\sum_{n\geq1}\frac1{n!}
   \sum_{\substack{d_1,\ldots,d_n\geq2\\
                   \sum_i(d_i-1)=2g-3+h}}
   \left(\int_{\MM_{g,n}}\lambda_{g-h}
                                  \prod_{i=1}^n\psi_i^{d_i}\right)
   \frac{\prod_{i=1}^n\mathcal I_{d_i}}
        {(1-\mathcal I_1)^{2g-2+n}}.
\end{equation}
The sum is finite. In genus one the positive-defect term is
\begin{equation}\label{eq:s07:genus-one-seed}
 \mathcal F_{1,1}=\mathcal G_{1,1}
                         =-\frac1{24}\log W_1.
\end{equation}
For every $g\geq h\geq1$, both
$\partial_x\mathcal F_{g,h}$ and
$\partial_x\mathcal G_{g,h}$ are rational expressions in the jets.
Every term has a denominator of degree at least two in $W_1$, and
vanishes when $W_k=0$ for all $k\geq2$.
\end{lmm}
\begin{prf}
For $h\geq1$ the string equation is homogeneous, by
Lemma~\ref{lmm:s07:string}. Its characteristic with parameter $a$
sends the times to
$\sum_{k\geq0}t_{d+k}a^k/k!-\mathbf1_{d=0}a$; here
$\mathbf1_{d=0}$ is one for $d=0$ and zero otherwise.
At $a=v$, its zero-degree time is zero by
\eqref{eq:s07:genus-zero-inverse}, and its remaining times are
$\mathcal I_d$. Hence $\mathcal F_{g,h}$ can be evaluated at these
times with no zero-degree insertions.

On a stable target, forgetting a degree-one marking gives the dilaton
factor $2g-2+n$. Geometrically, this is
$\pi_*\psi_{n+1}=(2g-2+n)[\MM_{g,n}]$; the cotangent pullback
corrections restrict to zero after multiplication by $\psi_{n+1}$
on the three-pointed rational tail. Summing $r\geq0$ such markings
with weight $\mathcal I_1^r/r!$ gives
$(1-\mathcal I_1)^{-(2g-2+n)}$. Dimension imposes
$\sum_i(d_i-1)=2g-3+h$ on the remaining markings. Each summand on
the left is at least one, so $n\leq2g-3+h$ and the sum is finite.
There is no $n=0$ term for $g\geq2$, since it would require
$g-h=3g-3$.

For $g=h=1$, dimension with no zero-degree insertions permits only
degree-one markings. On $\MM_{1,1}$, evaluation of a differential at
the marked point identifies the Hodge and cotangent lines, also on
the irreducible nodal fiber. Thus
$\int_{\MM_{1,1}}\psi_1=\beta_1=1/24$.
Starting with this one-point value, the dilaton equation gives
$\int_{\MM_{1,r}}\prod_{i=1}^r\psi_i=(r-1)!/24$ for $r\geq1$.
Its exponential generating function is
$-\log(1-\mathcal I_1)/24$, proving
\eqref{eq:s07:genus-one-seed}. The sine transformation has no
other term with $g=h=1$.

Finally, implicit differentiation gives
\begin{equation}\label{eq:s07:jet-x}
 \partial_xv=W_1^{-1},\qquad
 \partial_xW_k=\frac{W_{k+1}}{W_1}\quad(k\geq1).
\end{equation}
Applying a positive number of $x$ derivatives to
\eqref{eq:s07:finite-jets} produces rational expressions with the
stated denominators and at least one higher jet in every numerator.
The same holds for the logarithmic seed after its first derivative.
Since
\[
 \mathcal G_{g,h}
 =\sum_{k=0}^{g-h}\frac{(-1)^k}{2^{2k}(2k+1)!}
                           \partial_x^{2k}\mathcal F_{g-k,h},
\]
it also holds for $\partial_x\mathcal G_{g,h}$.
\end{prf}

Define, for $g\geq h\geq1$,
\begin{equation}\label{eq:s07:transport-definitions}
 L_{g,h}=\partial_x\mathcal G_{g,h},\qquad
 \mathscr D=\partial_{t_1}-v\partial_x.
\end{equation}
The notation $L_{g,h}(\theta;x)$ denotes specialization at
\eqref{eq:s07:times}. In particular,
$L_{1,1}=-W_2/(24W_1^2)$.
The identities for the transport operator on the jets are
\begin{equation}\label{eq:s07:jet-transport}
 \mathscr Dv=0,\qquad
 \mathscr DW_1=-1,\qquad \mathscr DW_k=0\quad(k\geq2).
\end{equation}
They follow by differentiating $W(v)=x$ in $t_1$ and using
$\partial_{t_1}W(u)=-u$. After specialization,
$\mathscr D=-\partial_\theta-v_\theta(x)\partial_x$.

We now integrate the field equation once in $x$. The finite jet
description fixes its possible constant, and its decay under
$W_1\mapsto W_1+\sigma$ fixes the subsequent characteristic integral.

\begin{prp}[Exact transport]\label{prp:s07:transport}
For every $g\geq1$,
\begin{equation}\label{eq:s07:transport-one}
 \mathscr DL_{g,1}=c_g\partial_x^{2g}v.
\end{equation}
For $g\geq h\geq2$,
\begin{equation}\label{eq:s07:transport-higher}
 \begin{aligned}
 \mathscr DL_{g,h}
 ={}&\sum_{k=1}^{g-h+1}c_k\partial_x^{2k+1}L_{g-k,h-1}\\
 &+\frac12\sum_{i=1}^{h-1}
   \sum_{\substack{g'+g''=g\\g'\geq i,\ g''\geq h-i}}
       (\partial_xL_{g',i})(\partial_xL_{g'',h-i}).
 \end{aligned}
\end{equation}
\end{prp}
\begin{prf}
Extract the coefficient of $z^g\eta^h$ in
\eqref{eq:s07:ilw}, using $U|_{\eta=0}=v$. The terms involving
$v$ give $\partial_x(\mathscr DL_{g,h})$ after the product rule.
The dispersive part gives the derivative of the linear term displayed
in \eqref{eq:s07:transport-one} or
\eqref{eq:s07:transport-higher}. The two positive-defect factors
in $U\partial_xU$ give the derivative of the ordered quadratic sum
with factor $1/2$. Thus the asserted equations hold up to an
$x$-independent expression.

That expression is a rational function of finitely many jets, by
Lemma~\ref{lmm:s07:finite-jets} and
\eqref{eq:s07:jet-transport}. The kernel of $\partial_x$ on this
ring consists of numerical constants. Indeed, in
\[
 \partial_x=\frac1{W_1}\sum_{k\geq1}W_{k+1}\partial_{W_k},
\]
the coefficient of the highest jet forces independence of the
preceding jet; descending induction leaves only a constant.
The jets are independent formal coordinates, as is seen at $x=0$
by varying $t_1,t_2,\ldots$. Setting all $W_k$ with $k\geq2$ to
zero makes both sides of the asserted equations zero, so the
constant is zero. This proves
\eqref{eq:s07:transport-one}--\eqref{eq:s07:transport-higher}.
\end{prf}

\begin{crl}[Characteristic recovery]\label{crl:s07:characteristic-recovery}
For integers $g\geq h\geq1$, real $\theta\geq0$ and $x$ in a sufficiently small complex
neighborhood of zero, set
$X_\sigma=x+\sigma v_\theta(x)$ for $\sigma\geq0$.
On the inverse branch obtained by keeping $v_\theta(x)$ fixed,
\begin{equation}\label{eq:s07:characteristic}
 v_{\theta+\sigma}(X_\sigma)=v_\theta(x),\qquad
 L_{g,h}(\theta+\sigma;X_\sigma)=O_{g,h,\theta,x}(\sigma^{-2}).
\end{equation}
The corresponding integral formulas are
\begin{equation}\label{eq:s07:characteristic-one}
 L_{g,1}(\theta;x)
 =c_g\int_0^\infty
       (\partial_x^{2g}v_{\theta+\sigma})(X_\sigma)\,\dd\sigma
\end{equation}
and, for $g\geq h\geq2$,
\begin{equation}\label{eq:s07:characteristic-higher}
 \begin{aligned}
 L_{g,h}(\theta;x)=\int_0^\infty
 \Biggl[&\sum_{k=1}^{g-h+1}c_k\partial_x^{2k+1}L_{g-k,h-1}\\
 &+\frac12\sum_{i=1}^{h-1}
   \sum_{\substack{g'+g''=g\\g'\geq i,\ g''\geq h-i}}
     (\partial_xL_{g',i})(\partial_xL_{g'',h-i})
 \Biggr](\theta+\sigma;X_\sigma)\,\dd\sigma.
 \end{aligned}
\end{equation}
All derivatives inside the brackets are taken before evaluation on
the characteristic. The integrals converge at each fixed $g,h$ and
also express identities of formal Taylor series at $x=\theta=0$.
\end{crl}
\begin{prf}
Along the stated curve, the inverse equation reads
$w_{\theta+\sigma}(v_\theta(x))=x+\sigma v_\theta(x)$.
The first jet becomes $w_\theta'(v_\theta(x))+\sigma$, and the
higher jets stay fixed. For $x$ sufficiently close to zero the
first jet has positive real part, so this defines the branch for
all $\sigma\geq0$. The finite jet expressions give the decay in
\eqref{eq:s07:characteristic}. The source in
\eqref{eq:s07:transport-one} is $O(\sigma^{-3})$, since
$2g\geq2$; the linear and quadratic sources in
\eqref{eq:s07:transport-higher} decay at least as fast.

The derivative along the characteristic is $-\mathscr D$.
Integrating from $\sigma=0$ to infinity and using the zero limit
therefore gives \eqref{eq:s07:characteristic-one} and
\eqref{eq:s07:characteristic-higher}. Each integrand is a finite
sum of rational functions of $W_1+\sigma$ with integrable powers.
The calculation is consequently also valid coefficientwise in the
formal jets, and in their specialized Taylor series.
\end{prf}

The first layer can be evaluated without a recursion. For an
auxiliary formal variable $y$, define the normalized logarithmic
difference quotient
\begin{equation}\label{eq:s07:logarithmic-kernel}
 \Lambda_\theta(x,y)
 =\log\frac{v_\theta(x+y)-v_\theta(x)}
                  {y\,\partial_xv_\theta(x)}.
\end{equation}
Its argument has constant term one. The normalization removes only
a term independent of $y$ and makes the formal logarithm unambiguous.

\begin{lmm}[The first-layer kernel]\label{lmm:s07:first-kernel}
For every $g\geq1$,
\begin{equation}\label{eq:s07:first-kernel}
 L_{g,1}(\theta;x)
 =c_g(2g-1)![y^{2g-1}]\Lambda_\theta(x,y).
\end{equation}
\end{lmm}
\begin{prf}
Fix $v=v_\theta(x)$, and use a local formal variable $u$ to put
$A(u)=(w_\theta(v+u)-w_\theta(v))/u$.
The constant term of $A$ is $W_1$. Lagrange inversion at
$X_\sigma$ gives
\[
 (\partial_x^{2g}v_{\theta+\sigma})(X_\sigma)
 =(2g-1)![u^{2g-1}](A(u)+\sigma)^{-2g}.
\]
Integration in \eqref{eq:s07:characteristic-one} yields
$c_g(2g-2)![u^{2g-1}]A(u)^{-(2g-1)}$.
If $V(y)=v_\theta(x+y)-v$, then $y=V A(V)$.
Lagrange--B\"urmann inversion applied to $-\log A(V)$ gives,
for each integer $N\geq1$,
\[
 [y^N]\log\frac{V(y)}y
                 =\frac1N[u^N]A(u)^{-N}.
\]
The same positive-degree coefficients occur in
\eqref{eq:s07:logarithmic-kernel}. Taking $N=2g-1$ proves the
claim. In particular $g=1$ gives $-W_2/(24W_1^2)$, in agreement
with the independently determined genus-one seed.
\end{prf}
